% !TEX root = main.tex


%%
%Railway infrastructure is troubled by defects occurring on the rail component of the track. Initiated cracks grow over time with the potential of breaking the rail. Rail breaks impose safety risks to rolling stock on the track in terms of increased chances of derailment. These risks lead railway infrastructure managers to enforce speed restrictions and/or track replacements with respect to the severity of the defect. Thus, rail defects are also associated with decreased network efficiency and an increased carbon footprint \citep{cannon2003a,Zoeteman2014DutchBeyond}. 
%To mitigate these risks, infrastructure managers classify defects into severity categories based on type and physical size \citep{UIC_2002_defecthandbook}. These classifications guide inspection intervals and maintenance interventions.
%
%Due to limited maintenance resources, railway infrastructure managers often face the challenge of deciding which defects to remove among a substantially larger collection of identified defects. Naturally, defects posing immediate safety risks in the near future should be prioritized. 
%However, this assessment can be difficult to make due to several reasons. 
%For one thing, as rail defects grow over time, the most recently recorded severity classification might not be accurate at the time of maintenance decision making. Secondly, defect growth is influenced by the local conditions of the railway network \citep{bai2017a}. For instance, high tonnage and heavy curvature, among other factors, accelerate defect development \citep{Kumar2008HolisticSweden}.  

\section{Introduction} \label{sec:introduction} %General background
Various infrastructure systems are critically affected by defects. Defects may grow over time and eventually cause fractures that make the infrastructure unsafe for operations.
For railway infrastructure, rail breaks pose serious safety concerns in terms of increased derailment risks, leading infrastructure managers to enforce speed restrictions or track replacements when advanced defects are discovered. 
Consequently, rail defects are also associated with reduced network efficiency and increased carbon footprint \citep{cannon2003a,Zoeteman2014DutchBeyond}.  

To mitigate these risks, infrastructure managers carry out inspections where substantial damage can be regarded as point defects. 
Defects are classified into severity categories based on type and physical size \citep{UIC_2002_defecthandbook}, which guide inspection intervals and maintenance interventions.  

Because maintenance resources are limited, managers often face the challenge of prioritizing reinspections and maintenance among a large number of identified defects.
Defects that pose imminent safety risks should be treated first, but this assessment is difficult for two reasons. 
First, since defects evolve over time, the last recorded severity classification may no longer be accurate at the time of decision making.
Second, defect growth is influenced by local operating conditions such as tonnage and curvature, which accelerate degradation \citep{bai2017a,Kumar2008HolisticSweden}.
These challenges motivate the development of models capable of forecasting the evolution of defect severity, explicitly accounting for uncertainty and defect-specific conditions.

We consider this a forecasting problem over an ordinal (and multicategory) state space representing severity classes. Since inspection data are interval-censored, and degradation rates vary with defect-specific conditions, the modeling problem involves predicting the probability of severity transitions under uncertainty, based on censored transition data and exogenous information.

Accurate forecasts of defect propagation that account for defect-specific conditions and inspection intervals can therefore inform timely maintenance decisions and improve both safety and efficiency \citep{Binder2023}. In this setting, a Markov jump process (MJP; also known as a continuous-time Markov chain), provides a dynamic probabilistic model: it yields full predictive distributions over future severity classes and can incorporate spatial variability through covariates \citep{bisgaard2024A}. Unlike point forecasts or expected-time-to-failure models, a full probabilistic forecast explicitly quantifies uncertainty, allowing infrastructure managers to balance risk and cost under resource constraints.

MJPs have been widely applied in infrastructure degradation modeling for bridges \citep{Mizutani2023RegimeSwitching}, pipelines \citep{JimenezRoa2024Nonhomogeneous}, and drainage systems \citep{Wu2021ConditionPrediction}, making them increasingly relevant for predictive maintenance. %Despite their advantages in forecasting, the predictive accuracy of CTMCs is rarely evaluated in railway applications. 

\subsection{Background} \label{sec:background}%Specific background
Statistical degradation models have been widely applied to rail infrastructure \citep{Davari2021AIndustry}.
\citet{ghofrani2019a} use a proportional hazards model to estimate the time to the reappearance of defects on a repaired track segment. The study lumps degradation into a binary classification problem, that is failure mode or not, by defining thresholds for the size and magnitude of the defects. Conceptually, this is equivalent to the approach of \citet{dick2003a} that apply logistic regression to estimate the probability of service failures for each rail segment. While interpretable, both approaches collapse defect growth into binary outcomes and do not model the propagation of individual defects across severity classes, nor are they applied for forecasting.

Other approaches quantify risk or capture limited forms of progression, but do not fully represent the defect dynamics observed in rail inspection data. For example, \citet{jamshidi2016a} estimate the failure risk of surface defects using Bayesian regression, but do not model the sequential or ordered progression between intermediate severity states. \citet{bressi2021a} apply discrete-time Markov chains to ballast degradation, which do model progression but are poorly suited to the temporally irregular inspection intervals common in rail defect tracking.

Markov jump processes have been proposed for modeling severity progression in rail defects \citep{podofillini2006a} and bridge deterioration \citep{sun2024a}, assuming constant transition rates and transitions only between adjacent states. However, these models do not include covariates, limiting their ability to capture spatial variability. \citet{bisgaard2024A} extends this framework to incorporate covariate-dependent transition rates and interval-censored observations, using a generalized Erlang structure of the generator matrix and maximum likelihood estimation.

While these studies establish important modeling foundations for degradation processes, they largely focus on parameter estimation and simulation. None explicitly assess the forecast accuracy of their models, despite the increasing focus on data-driven predictive maintenance \citep{Binder2023}. Forecast evaluation requires the use of proper scoring rules \citep{Gneiting2007} and well-chosen benchmarks \citep{Hyndman2014}. Probabilistic scoring rules such as the logarithmic score and ranked probability score are standard tools for evaluating ordinal or multicategory forecasts \citep{Epstein1969,friedman1983a}. However, benchmarking multicategory forecasts remains underdeveloped in the context of railway degradation, particularly given the lack of suitable baseline models.

Beyond railway applications, MJPs have been widely used for predictive modeling in other infrastructure domains. \citet{Mizutani2023RegimeSwitching} propose a regime-switching MJP for bridge deterioration, improving fit by allowing transitions between distinct aging regimes. \citet{JimenezRoa2024Nonhomogeneous} assess time-inhomogeneous MJPs for sewer pipe forecasting and show that non-constant transition rates outperform homogeneous models in predictive accuracy. \citet{Wu2021ConditionPrediction} apply MJPs to drainage assets and demonstrate that including covariates improves forecasts. %These studies underscore the importance of model structure in CTMC-based prediction.

Despite this progress, several aspects remain underexplored in infrastructure degradation modeling.
Systematic evaluations of alternative generator parameterizations (such as whether to allow skip transitions or restrict jumps to neighboring states) remain scarce, with most studies adopting a single structure without benchmarking predictive performance against alternatives, e.g. \citep{bisgaard2024A,Zbiciak2025}.
Second, while covariate-augmented MJPs have been proposed in related fields \citep{Wu2021ConditionPrediction}, the forecasting impact of covariates has not been systematically quantified in railway degradation contexts.
Third, estimation is typically performed by maximum likelihood or least squares, whereas scoring-rule-consistent estimation (optimum score estimation) is rarely used, even though it is known to align estimation with evaluation and improve forecast performance \citep{Gneiting2007,friedman1983a,Copas1983}.
Fourth, predictive assessment has relied mainly on point-based evaluation \citep{JimenezRoa2024Nonhomogeneous,Mizutani2023RegimeSwitching}, while probabilistic scoring rules (e.g., log score, ranked probability score) provide a means to assess both the accuracy and robustness of forecasts.

%Finally, the idea of a state-dependent time scaling remains largely absent from degradation modeling. In this construction, the effective “clock” of the process runs faster or slower depending on the current defect state \citep{AITSAHALIA20101}, allowing early stages of crack growth to evolve slowly and advanced stages to progress more rapidly. This mechanism does not replace the exponential assumption of CTMC holding times, but rescales them state-dependently, thereby offering additional flexibility to represent heterogeneous degradation speeds across severity classes. 
Some prognostics literature raise the concern that classical Markov models often fail to capture realistic sojourn behavior across states \citep{AIZPURUA2022}. Related work on time-varying operating conditions has explored alternative ways of rescaling degradation rates through covariates \citep{liu2024a,bisgaard2024A}, but to our knowledge no prior studies within reliability of infrastructure have modified the sojourn time estimates within the Markov framework to account for heterogeneity across degradation states, nor examined how such adjustments influence probabilistic forecast accuracy.



% cite Zbiciak2025, for not evaluating predictive accuacy (else it is used to predicting Service Life of Buildings and Structural Components)


%Beyond railway applications, continuous time Markov chains have also been used for predictive modeling in other infrastructure domains. \citet{Mizutani2023RegimeSwitching} propose a regime-switching CTMC for bridge deterioration, improving fit by allowing transitions between distinct aging regimes. \citet{JimenezRoa2024Nonhomogeneous} assess time-inhomogeneous CTMCs for sewer pipe forecasting and show that non-constant transition rates outperform homogeneous models in predictive accuracy. \citet{Wu2021ConditionPrediction} apply CTMCs to drainage assets and demonstrate that including covariates improves forecast accuracy. These studies underscore the importance of model structure in CTMC-based prediction.
%
%Despite this progress, a systematic evaluation of CTMC parameterization choices—such as whether to allow skip transitions or include covariates—remains largely unexplored. Most infrastructure studies do not benchmark different transition structures or analyze how model assumptions affect forecast accuracy. Moreover, forecast models are rarely estimated with scoring-rule-consistent objectives, even though optimum score estimation is known to improve predictive performance by aligning estimation with the chosen evaluation metric \citep{Gneiting2007,friedman1983a,Copas1983}. While \citet{JimenezRoa2024Nonhomogeneous} compare homogeneous and nonhomogeneous CTMCs using RMSE, and \citet{Mizutani2023RegimeSwitching} demonstrate improved fit via switching regimes, none apply proper scoring rules or investigate how alternative generator structures impact multicategorical forecast performance.



%\citet{bressi2021a} apply discrete-time finite state Markov chains to model defect propagation in railway ballast. This model formulation has the ability to model state-dependency as well as temporal dependency. However, the discrete-time formulation of the Markov chains would cause issues in adapting the framework to rails since measurements of rail defects are recorded temporally irregular.
%\citet{podofillini2006a} propose a continuous-time Markovian framework to model the propagation of rail defects.
%The study assumes a single absorbing state corresponding to complete rail failure, while the non-absorbing states represent defect classes ordered by severity.
%The model is used to determine optimal measurement policies, but on generic data. \citet{sun2024a} apply a similar model approach, but for the deterioration of bridge defects. Although the Markovian approaches of \citet{podofillini2006a} and \citet{sun2024a} are capable of describing the probabilistic interdependency among the defect categories, they do not take exogenous information into account, why the approaches would not be capable of capturing the spatial heterogeneity of defect propagation in a railway network.
%\citet{bisgaard2024A} extends the Markovian model frameworks of \citet{podofillini2006a} and \citet{sun2024a} to model defect propagation subject to exogenous dependency, hereby letting the transition rate be a function of the properties of the rail component at its specific location in the network. As the timing of the jumps are unknown, the study considers interval-censored data (or partially observable defect trajectories).
%Exogenous dependency (i.e. intrinsic rail characteristics like the ones included in \citet{ghofrani2019a} or \citet{dick2003a}) are introduced to the transition rates, in a fashion similar to the parameterizations of conditional probabilities in \citet{bisgaard2024B}. 
%While \citet{podofillini2006a} and \citet{sun2024a} focus on determining optimal measuring policies, \citet{bisgaard2024A} focus on obtaining robust estimates for the transition rate matrix, given industrial data.
%Both \citet{podofillini2006a}, \citet{sun2024a}, and \citet{bisgaard2024A} assume the generalized Erlang distribution with distinct transition rates for each jump in the state space, with jumps only allowed between neighboring defect classes. Neither of the studies, consider different parameterizations for the generator.
%\citet{bisgaard2024A} argues that the proposed Markovian framework is suitable for forecasting.
%Neither
%\citet{bressi2021a}, \citet{podofillini2006a}, \citet{sun2024a}, or \citet{bisgaard2024A} apply the proposed models for prediction.
%To realize predictive maintenance for the railway, it is not only important to develop models; it is also important to evaluate their forecast performance \citep{Binder2023}. 
%Proper assessment of a forecast model involves at least two vital components, i.e. 1) choosing a suitable metric, and 2) choosing suitable reference predictors \citep{Hyndman2014}. 
%As the Markov jump process is of probabilistic type, it predicts distributions, which in particular is useful for robust maintenance decision-making via the forecast's inherent uncertainty quantification.
%Therefore, probabilistic methods should also be used for evaluation \citep{Gneiting2007}. 
%The ordinal and multicategorial state space of the Markov jump process gives rise to further considerations in terms of forecast evaluation.
%Multiple (and widely discussed) metrics exist for such purposes, e.g. the Brier Score, Log Score, and Ranked Probability Score, among others \citep{Brier1950,Epstein1969,friedman1983a}. 
%Good-practices for benchmarking continuous state space forecasts are well-established, such as the persistence model, random walk, or ARIMA \citep{forecast_evaluation_IJF,hewamalage2022forecast}. 
%However, it is less clear which reference predictors are suitable to use to benchmark multicategorial probability forecasts. 
%In addition to the interest of practitioners within probabilistic forecasting and applied matrix exponential distributions, such an investigation is relevant for railway infrastructure managers because the evaluation of the model's predictive capabilities is needed to realize well-functioning predictive maintenance \citep{Binder2023}.


%\citet{jamshidi2016a} propose a defect-based risk assessment for rail failures, which is quantifying the probability that a surface defect will develop into a rail break. The study divides surface defects into four categories based on size and calculates the rail failure risk using Bayesian inference of an exponential decay regression type of model. Although \citet{jamshidi2016a} consider multiple defect categories modeling how each category relates to the failure mode, the model formulation does not capture how the categories relate to each other. In that way, the model ignores the temporal dependency between the defect categories, and that mildly severe defects can develop into semi-severe categories, before it reaches the total failure mode of rail break. The study does not consider defect prediction.


\subsection{Contribution}
Reliable forecasting of defect severity is essential for prioritizing maintenance under constrained resources. Since severity classifications guide intervention decisions, models that accurately predict how defects will progress can help infrastructure managers decide which defects to reinspect or repair first. However, without systematic evaluation, it is unclear how different modeling assumptions (such as transition structure, state-dependent time scaling, regression on exogenous information, and mixture distributions) impact predictive performance. To our knowledge, there has been little systematic effort to evaluate predictive performance in probabilistic railway degradation modeling. 
The main contributions of this paper are:

\begin{itemize}
    \item Generator structure: We extend prior MJP-based degradation models by testing alternative parameterizations of the generator.    
    \item Modified sojourn times: A heuristic extension of the Markov jump process that modifies the sojourn time estimates is introduced. This allows more flexibility in representing deterioration dynamics across severity classes.
      \item Mixture model: A mixture MJP comprised of latent components with corresponding generators is developed to model complex transition behavior in the data.
    \item Exogenous dependencies: We test whether adding local track covariates improves forecasts, comparing a global-coefficient regression to a state-dependent-coefficient regression.
     \item Optimum score estimation: Model estimation is aligned with the chosen forecast scoring rule. This approach improves consistency between model estimation and how predictions are evaluated.
     \item Forecast evaluation: We perform scoring-rule-based evaluation of the forecasts, benchmarking the Markov jump process across model specifications. Calibration, discrimination, and sharpness are also assessed. %To our knowledge, no prior degradation studies have conducted such an evaluation.
%    \item Implementation and reproducibility: An open-source implementation of the model and scoring framework in \texttt{R}/\texttt{C++} is provided, including three methods for computing the transient distribution (uniformization, Padé approximation, and eigendecomposition).
    \item Stochastic optimization program: We demonstrate how the predictive distributions can be embedded in a stochastic optimization program that plans inspection and maintenance under budget constraints. This connects forecasting to decision-making for infrastructure managers.
\end{itemize}

By combining statistical modeling, predictive evaluation, and decision-oriented output, the study contributes a robust methodology for condition-based forecasting in railway maintenance. The approach is broadly applicable to other infrastructure systems where asset condition is recorded in severity categories and inspection intervals are irregular.



% In this study, we build upon the proposed covariate-dependent Markov jump process evaluating its performance in predicting the propagation of rail defect severity. This study is the first of its kind to conduct such an investigation. We perform probabilistic forecast evaluation based on three different scoring rules and compare with multiple reference predictors consisting of simple methods and static regression models. Furthermore, we provide a diagnostic analysis of forecast quality investigating calibration and sharpness. 
% The proposed evaluation framework constitutes a rigorous approach in assessing the prediction quality of a proposed forecast model.
% We extend the proposed Markovian model framework of \citet{podofillini2006a}, \citet{sun2024a}, and \citet{bisgaard2024A} by reparameterizing the transition matrix allowing larger jumps in the state space than assumed by the generalized Erlang distribution. We evaluate the predictive accuracy of the various parameterizations and investigate the implications of the model assumptions on the transient distributions of the Markov chains. We have not found other studies conducting such predictive investigation on the influence of the generator parameterization.
% Furthermore, we quantify the gains in predictive accuracy by taking local conditions of the rail network into account. 
% The results of this study are relevant for railway infrastructure managers as accurate predictions on the propagation of rail defect severity increases safety and efficiency of railway track. The methodology has the potential for application in other domains experiencing structural asset degradation which is recorded by classification procedures.
% Lastly, this study contributes with a code repository where score computations, estimation methodology, and a prediction and evaluation framework are implemented for the Markov jump process and the reference predictors. The parameters of the Markov jump processes are determined using optimum score estimation based on the user specified combination of scoring rules. The implemented calculation method of the transient distribution relies on three different variations, i.e. uniformization \citet{sherlock2021a}, Padé approximations, and eigendecompositions of the generator \citep{bisgaard2024A}. The model implementations are highly efficient and can be adapted to other empirical settings with little \texttt{R} and \texttt{C++} experience.



\subsection{Paper outline}
The rest of this article is outlined as follows. In \Cref{sec:data}, we describe the empirical setting of the study, i.e. the recorded defect trajectories as well as included exogenous rail characteristics. We further qualitatively outline the prediction task addressed in the study. In \Cref{sec:methods}, we describe the mathematical model settings. This description includes the formulation of the Markov jump process in \Cref{sec:mjp} with its various generator parameterizations in \Cref{sec:generator_params}, estimation methodology in \Cref{sec:optimum_score_estimation} and \Cref{sec:error_metrics}, and calculation methods for the transient distributions in \Cref{sec:computation_remarks}. In \Cref{sec:reference_forecasts}, we present all the reference predictors. In \Cref{sec:prediction_error}, we estimate the prediction errors of the models wrt. the different scoring rules using a cross-validation procedure. The implications of the generator parameterizations are discussed in \Cref{sec:results_transient_dist_markov}. In \Cref{sec:reliability_analysis}, we provide a diagnostic analysis of the sharpness and calibration of the best performing prediction models. \Cref{sec:proposition_for_use} outlines a stochastic program taking the predictive distributions as input. Discussion and conclusions follow in \Cref{sec:conclusion}. 